\textbf{For practitioners.} The shortcut is safe only when the admissible embedding is unique, and on the third-party queries it rarely is (Section~\ref{sec:firstmatch}). A cheap safeguard follows from the reduction: after a first match, issue the counterexample query. If it completes without a witness, the first-match answer was right; if it finds one, the answer is \Unk{}; if it times out, report the result as undetermined rather than as the first-match label.

\textbf{For solver designers.} The suites reward different engines: LAD's propagation on large unanchored ambiguous patterns, Glasgow's staged preprocessing on the largest hosts, and SAT and cautious reasoning when label domains are narrow or the identity space is small enough to ground; whole-model CP wins some anchored strata. IdCon-Hard is a short-budget suite: with a minute per task the benchmarked methods jointly leave \hardVBSAllUnsolvedSixty{} of its 504 tasks open, and unpublished methods already do better at 2\,s (Appendix Table~\ref{tab:legacy}). IdCon-Hard-X moves the separating budget up: within 10\,s the best single method answers \hardxBestSingleTenN{} of its \hardxComplete{} assessed tasks and all methods together leave \hardxVBSUnsolvedTen{} open, while \hardxBestSingleSixty{} answers \hardxBestSingleSixtyN{} with 60\,s. IdCon-Cite is a 2\,s suite: some method answers every task within 10\,s, but at 2\,s the best single method answers \citeBestSingleTwoN{} of \citeTasks{}. None of the suites is beyond exact search with a minute per task, so IdCon measures the cost of exact consensus rather than its feasibility; the released generator scales further. The hardest outcomes are the universal ones, where proving \Sat{} requires refuting every non-$S$ placement; the independent check shows that DRAT proofs for such refutations are practical for many IdCon instances, so proof-producing consensus solvers \cite{auditablecp} can be held to certified negatives.
